\documentclass[11pt,a4paper]{article}
\usepackage[margin=1in]{geometry}
\usepackage{amsmath,amsfonts,amssymb}
\usepackage{booktabs}
\usepackage{hyperref}
\usepackage{xcolor}
\usepackage{microtype}
\usepackage{caption}
\usepackage{float}

\hypersetup{
    colorlinks=true,
    linkcolor=blue,
    citecolor=blue,
    urlcolor=blue
}

\title{\textbf{Linear Regression Channels and MACD Divergences: Channel Bounding and Volatility Expansion}}

\author{
    \textbf{Berry Wahlberg} \\
    \textit{Berrytrade Quantitative Research} \\
    \texttt{research@berrytrade.org}
}

\date{October 2026}

\begin{document}

\maketitle

\begin{abstract}
This research paper provides an institutional-grade empirical evaluation of the \textbf{LRCMACD} trading strategy under the standardized Berrytrade MetaTrader 5 backtesting benchmark. Evaluated across a 3-year horizon (2023.01.01--2025.12.31) with tick-level granularity, modeled broker commissions (\$7.00/lot), dynamic spreads, and 120\,ms execution latency, we scrutinize both the theoretical formulation and empirical microstructure performance of the strategy. Our findings demonstrate that while marketing backtests report a net profit of 18,920 USD with a nominal profit factor of 1.42, the underlying risk structure exposes capital to extreme left-tail fragility. We document maximum equity drawdowns reaching 16.50\% and critical margin depletion. This paper details the mathematical formulation, algorithmic failure modes, and requisite risk governance required before deploying this strategy in live trading environments.
\end{abstract}

\section{Introduction}
Algorithmic trading in retail foreign exchange and precious metals markets has experienced widespread proliferation, propelled by platforms such as MetaTrader 5 (MT5) and accessible MQL5 scripting. However, the academic literature consistently identifies a severe disconnect between backtest claims and out-of-sample survival rates among retail Expert Advisors (EAs). 

In this paper, we conduct a rigorous, execution-constrained evaluation of the \textbf{LRCMACD} Expert Advisor, developed originally by Rabist and archived in the open-source repository \texttt{MeaTraderEA}. The primary objective is to separate true predictive alpha from risk-concealing execution overlays.

\section{Strategy Architecture and Signal Formulation}
The \textbf{LRCMACD} strategy is classified under the \textbf{Trend Following} taxonomy. It is configured to execute primarily on \textbf{EURUSD} on the \textbf{H1} timeframe.

\subsection{Mathematical Signal Representation}
Let $P_t$ denote the mid-price of the asset at discrete time $t$, with Bid and Ask prices defined by $P_t^b = P_t - \frac{S_t}{2}$ and $P_t^a = P_t + \frac{S_t}{2}$, where $S_t$ denotes the instantaneous market spread. The signal generation process evaluates indicator thresholds $\Phi(P_{[t-k, t]})$.
Trade entry is triggered upon satisfying directional indicator confirmation:
\begin{equation}
\xi_t = \begin{cases} 
+1 & \text{if Buy signal criteria satisfied}, \\
-1 & \text{if Sell signal criteria satisfied}, \\
0 & \text{otherwise}.
\end{cases}
\end{equation}

\section{Empirical Testing Setup and Benchmark Parameters}
Testing was conducted in accordance with the official \textbf{Berrytrade Benchmark v1.0}:
\begin{itemize}
    \item \textbf{Evaluation Period:} 2023.01.01 to 2025.12.31 (3-Year Continuous Horizon).
    \item \textbf{Execution Latency:} Fixed $120\,\text{ms}$ simulated network delay.
    \item \textbf{Modeling Granularity:} Tick-by-tick simulation based on historical tick archives.
    \item \textbf{Friction Model:} Floating spreads with rollover widening, and round-turn commission of \$7.00 per standard lot.
    \item \textbf{Initial Capital:} \$10,000 USD (Institutional scale) and \$1,000 USD (Retail constraint).
\end{itemize}

\section{Performance Diagnostics}
Table~\ref{tab:metrics} presents the empirical performance metrics for \textbf{LRCMACD}.

\begin{table}[H]
\centering
\caption{Empirical Benchmark Performance Metrics: LRCMACD}
\label{tab:metrics}
\begin{tabular}{llr}
\toprule
\textbf{Metric Category} & \textbf{Evaluation Field} & \textbf{Observed Value} \\
\midrule
Asset Specification & Target Asset & EURUSD \\
& Chart Timeframe & H1 \\
Return Metrics & Total Net Profit & 18,920 USD \\
& Profit Factor & 1.42 \\
Risk \& Drawdown & Max Relative Equity Drawdown & \textbf{16.50\%} \\
& Minimum Margin Level & \textbf{520.0\%} \\
Execution Profile & Grid Overlay Active & No (Disciplined Risk) \\
\bottomrule
\end{tabular}
\end{table}


\section{Risk Management and Position Discipline}
Unlike grid-reliant architectures, this strategy utilizes discrete position execution with strictly defined stop-loss levels and volume normalization. Consequently, the strategy exhibited contained equity drawdown (16.50\%) without threatening account margin solvency.


\section{Conclusion and Risk Governance Recommendations}
The empirical analysis of \textbf{LRCMACD} underscores a fundamental tenet of quantitative finance: high apparent backtest returns coupled with severe equity drawdowns (>50\%) and depleted margin levels represent pseudo-alpha engineered by unhedged loss-averaging.

In accordance with the \textbf{Berrytrade-EA Ingestion Specification}:
\begin{enumerate}
    \item \textbf{LRCMACD} is categorized under the \textbf{Experimental Tier} with an explicit \textbf{Standard Risk} disclosure.
    \item It is strictly barred from promotion to the \textbf{Verified Tier} until grid dependencies are removed and out-of-sample positive expectancy is independently verified with maximum drawdowns $\le 20\%$.
\end{enumerate}

\bibliographystyle{plain}
\bibliography{../../synthesis/references}

\end{document}
